% %==============================================================================
% %  CHAPTER 4 — DATA, CALIBRATION AND VALIDATION (methodology)
% %------------------------------------------------------------------------------
% %  Style: matches introduction.tex / model.tex (\citep/\citet authoryear; keys
% %  resolve against introduction_refs.bib + model_refs.bib). Written while the
% %  final calibration runs are now LOCKED (D65, output/relock/): this chapter fixes the
% %  METHODOLOGY; calibrated values, fit tables, loss surfaces, and the ablation/robustness
% %  numbers are reported in the results chapter. [NUMBERS: ...] marks the spots
% %  to fill from output/relock/ (and output/relock_gbm/ for the SV-MJD robustness).
% %==============================================================================

This section presents the methodology for calibrating the agent based financial market simulation. As this serves as the first stage and base of the clearing layer, the market simulation should be well calibrated and able to reproduce the features of real market returns. A common method to do this, and the one adopted by this model, is the reproduction of the moments of empirical returns through the simulation. This chapter describes the data used, the parameters calibrated, the method of calibration and the calibration, and a validation of the calibrated parameters. The chosen moments are the stylised facts of returns we see empirically - the distribution of returns, heavy tails, absence of autocorrelation and volatility clustering. 

\section{Data and Regimes}

The dataset used for calibration is the E-mini S\&P futures at a one minute frequency, matching the model time steps. There are two separate regimes, a calm regime and a stressed regime. The calm regime spans the a calm stretch in late 2019, and the stressed regime includes data from the early Covid-19 window, which was referred to as the most significant test of the resilience of financial markets since the GFC \citep{bcbs_cpmi_iosco_2022}. Parameters are calibrated separately on both regimes, as they differ substantially in volatility, microstructure and return distribution. Returns are intraday log-returns of the mid price, and some summary statistics are reported in \ref{tab:cal-data}.

\begin{table}[h] 
\centering
\small
\setlength{\tabcolsep}{5pt}
\caption{\textbf{Empirical summary statistics for the two calibration regimes.}}
\label{tab:cal-data}
\begin{threeparttable}
\begin{tabular}{lccccccccc}
\toprule
 & & \multicolumn{4}{c}{One-minute returns} & \multicolumn{4}{c}{Daily returns}\\
\cmidrule(lr){3-6}\cmidrule(lr){7-10}
Regime & Sessions & Mean & SD & Skew. & Ex.\ kurt. & Mean & SD & Skew. & Ex.\ kurt.\\
 & & (bps) & (bps) & & & (\%) & (\%) & &\\
\midrule
Calm (2019) & 75 & 0.00 & 2.40 & $-1.41$ & 55.0 & 0.10 & 0.60 & $-0.85$ & 1.66 \\
Stressed (2020) & 73 & $-0.03$ & 12.40 & $-1.24$ & 67.1 & $-0.13$ & 3.14 & $-0.31$ & 2.50 \\
\bottomrule
\end{tabular}
\begin{tablenotes}[flushleft]\footnotesize
\item One-minute returns exclude overnight gaps; daily returns are close-to-close. bps$=10^{-4}$.
\item Both regimes hold $\approx$29{,}000 one-minute observations (29{,}961 vs 28{,}866) over a
comparable number of full RTH sessions (75 vs 73), so the two samples are of matched size.
\end{tablenotes}
\end{threeparttable}
\end{table}

Limit order depth placement follows a geometric law calibrated besides the other parameters, but this can also be directly estimated from L2 data. 

\section{Loss Function: Stylised Facts Distance}
\label{sec:cal-loss}

The market layer is calibrated to reproduce the stylised facts of returns we see in markets \citep{cont_2001} such as heavy tails, near-zero autocorrelation, and volatility clustering. Five different moments are used as proxies for these stylised facts, and are computed on simulated and empirical one minute log-returns. Parameters are chosen to reduce the distance between a weighted loss function defining these stylised facts in the simulation and in empirical data, similar to \cite{gao_2022_xgb}. The function there is further extended by the Hill Tail index used in \cite{gao_2022_hfabm}, and the weighted loss function takes the form:
\begin{equation}
\label{eq:cal-D}
D(\theta) = w_{\mathrm{KS}}\,\mathrm{KS}(\theta) + w_{V}\,\Delta V(\theta)
          + w_{\mathrm{ACF1}}\,\Delta\mathrm{ACF1}(\theta)
          + w_{\mathrm{ACF2}}\,\Delta\mathrm{ACF2}(\theta)
          + w_{\mathrm{Hill}}\,\Delta\mathrm{Hill}(\theta).
\end{equation}
The calibrated parameter vector is the minimisation of this distance over the admissible set $\Theta$,
\begin{equation}
\label{eq:cal-argmin}
\hat\theta = \arg\min_{\theta\in
\Theta} D(\theta).
\end{equation} 

\noindent A description of the five stylised facts are given below.

\paragraph{Distribution (KS).} The Kolmogorov--Smirnov statistic, from the KS test which tests for the equality of probability distributions.
$\mathrm{KS}=\sup_x|F_{\mathrm{sim}}(x)-F_{\mathrm{emp}}(x)|$ is the largest gap (supremum) between the simulated
and empirical return distribution functions. Matching it forces the simulated returns to share the
distribution function of the empirical distribution, which is an estimate of the cumulative distribution function \citep{gao_2022_xgb}.

\paragraph{Volatility ($\Delta V$).} The absolute difference in return standard deviation between the empirical and simulated returns
\begin{equation}
\Delta V = \big| \sigma_{\mathrm{sim}} - \sigma_{\mathrm{emp}} \big|,
\end{equation}
where $\sigma$ denotes the standard deviation of the one-minute return series.

\paragraph{Return autocorrelation ($\Delta\mathrm{ACF1}$).} Empirical returns have near-zero autocorrelation. To reproduce this, we take the mean absolute difference between the simulated and empirical return autocorrelations over the lags $L=\{1,5,10,20\}$ minutes,
\begin{equation}
\Delta\mathrm{ACF1} = \frac{\sum_{\ell\in L}\big|\,\mathrm{ACF}_{sim}(\ell,r)-\mathrm{ACF}_{emp}(\ell,r)\,\big|}{|L|}.
\end{equation}
Following \citet{gao_2022_xgb}, each $\mathrm{ACF}(\ell,r)$ is smoothed over three adjacent lags,
$\mathrm{ACF}(\ell,r)=\tfrac13(a_\ell+a_{\ell+1}+a_{\ell+2})$ with $a_k$ the lag-$k$ return
autocorrelation, so lags $\{1,2,3,\,5,6,7,\,10,11,12,\,20,21,22\}$ enter the function. The
component penalises predictability such as a bid-ask bounce.


\paragraph{Volatility clustering ($\Delta\mathrm{ACF2}$).} In empirical data, we see that volatility clusters over time, large returns are followed by large returns and low returns by low returns. This can be measure by studying the autocorrelation of absolute or squared returns. Autocorrelations of absolute returns are positive and decay slowly. $\Delta\mathrm{ACF2}$ is the same distance with the autocorrelations $\mathrm{ACF}(\ell,|r|)$
taken on $|r|$ (same lags, same three-lag smoothing),
\begin{equation}
\Delta\mathrm{ACF2} = \frac{\sum_{\ell\in L}\big|\,\mathrm{ACF}_{sim}(\ell,|r|)-\mathrm{ACF}_{emp}(\ell,|r|)\,\big|}{|L|}.
\end{equation}
For the dataset, squared returns were dominated by extreme values making their autocorrelation noisy, and hence absolute returns are used

\paragraph{Heavy Tails (Hill).} Tail heaviness is measured by the \citet{hill_1975} estimator of the
tail index of absolute returns, $\mathrm{Hill}$ (a lower value means a fatter tail), and its absolute difference is taken,
\begin{equation}
\Delta\mathrm{Hill} = \big| \mathrm{Hill}_{\mathrm{sim}} - \mathrm{Hill}_{\mathrm{emp}} \big|.
\end{equation}
The estimate is averaged over several tail cut-offs for robustness. \\

This combination of the KS statistic and Hill tail index is deliberate, using just the KS statistic makes tails very large and using just the Hill Tail index makes kurtosis unrealistically high. This is a deviation from \cite{gao_2022_hfabm} that use only the Hill tail index for their high-frequency LOB and \cite{gao_2022_xgb} that use only the KS statistic, and is motivated by the more granular step size alongside the importance of realistic tails due to the study of stressed periods.

\paragraph{Weighting.} The five components live on different scales --- a dimensionless KS statistic,
a standard deviation of order $10^{-4}$, autocorrelations, a tail index near three --- so the raw sum
in \eqref{eq:cal-D} would be dominated by the largest-scale term. Following
\citet{franke_westerhoff_2012}, each weight is the reciprocal of that component's own \emph{sampling}
standard deviation,
\begin{equation}
\label{eq:cal-w}
w_c = \frac{1}{s_c}, \qquad c \in \{\mathrm{KS},\,V,\,\mathrm{ACF1},\,\mathrm{ACF2},\,\mathrm{Hill}\},
\end{equation}
so every term in \eqref{eq:cal-D} is measured in units of sampling noise: an intrinsically noisy
moment is down-weighted and a precisely-measured one counts for more. The $s_c$ come from a
moving-block bootstrap of the empirical returns \citep{kunsch_1989} --- the series is resampled in
day-long blocks (long enough to preserve the dependence the autocorrelations measure), the whole
battery is recomputed on each resample, and $s_c$ is its standard deviation across resamples --- and
are fixed once per regime, so the loss is stationary and comparable across evaluations.
\citet{franke_westerhoff_2012} use the full inverse-covariance matrix and \citet{gao_2022_hfabm} the
inverse \emph{variance} $1/s_c^2$; the latter is scale-pathological here (the minute return-SD's tiny
sampling variance would swamp the sum), so the inverse standard deviation is the stable middle choice.
A value of $D$ near one is, loosely, a one-standard-error miss per moment.

